\documentclass[10pt,a4paper]{amsart}
\usepackage[margin=19mm]{geometry}
\usepackage{amsmath,amssymb,mathtools}
\usepackage[scr=rsfs,cal=euler]{mathalfa}
\usepackage{xfrac}
\usepackage[unicode,hidelinks,bookmarks=false]{hyperref}
\newcommand{\ie}{i.e.\ }
\newcommand{\cf}{cf.\ }
\newcommand{\resp}{respectively\ }
\newcommand{\Subsection}{\subsection}
\providecommand{\nobreakdash}{\nobreak\hbox{-}}
\setlength{\emergencystretch}{2em}
\multlinegap0pt
\usepackage[shortlabels]{enumitem}
\usepackage{amssymb}
\usepackage{xcolor}
\usepackage{tikz}
\usepackage[normalem]{ulem}
\newtheorem{lemma}{Lemma}
\newcommand{\F}{\mathbb{F}}
\renewcommand{\P}{\mathbb{P}}
\newcommand{\Span}[1]{\langle #1 \rangle}
\newcommand{\qbound}[1]{{\color{cyan}$q > {#1}$}}
\newcommand{\sm}{\mathrm{sm}}
\title{Effective Bertini theorems and zeros of $p$-adic forms of degrees 7 and 11}
\author{Lea Beneish, Christopher Keyes}
\date{Source: arXiv:2508.20192v2 (CC BY 4.0)}
\begin{document}
\maketitle

\noindent\emph{Source and licence.} Effective Bertini theorems and zeros of $p$-adic forms of degrees 7 and 11. Version arXiv:2508.20192v2 (CC BY 4.0), distributed by arXiv under the Creative Commons Attribution 4.0 International licence, http://creativecommons.org/licenses/by/4.0/, as recorded in the arXiv licence metadata for this exact version and shown on the abstract page https://arxiv.org/abs/2508.20192v2. Full licence notice: \texttt{licenses/arxiv-2508.20192-CC-BY-4.0.txt}.\par\smallskip

\noindent\textit{Editorial scope.} Counted unit: the article's lemma lem:plane\_with\_3pts\_d=11 (source0.tex lines 640--643). The article's complete original proof of it is delivered with it (source0.tex lines 645--691, one \texttt{proof} environment of 47 lines). No mathematics is written, repaired or re-derived: the statement and the proof are the article's own lines, in the article's own notation, and no retained line is substituted or re-flowed.  Retained uncounted as the source-local context the proof needs: lines 485--491; lines 550--553.  Nothing was omitted from any retained span, so the package carries no omission record.  No retained line contains a citation, so the wrapper carries no bibliography and no bibliography entry.  No retained line contains a figure environment, an image-inclusion command or drawing code, so no image is delivered and the package declares no asset dependency.  Editorial formatting only, in addition: an AMS \texttt{amsart} preamble that restates the article's own declarations and macros used by the retained lines, namely the environments \texttt{lemma} on separate counters, together with the standard packages those lines need.  The retained environments print in this document as Lemma 1 in order of appearance, whereas the article prints the same items with the numbering of its own full document.  The complete native source of the article as distributed in the arXiv e-print for this exact version is bundled unchanged as \texttt{sources/184007/source0.tex}.
\begin{lemma}
\label{lem:finding_sm_pt}
    Let $C \subset \P^2$ be a plane curve defined over $\F_q$ of degree $d$. Suppose $C' \subset C$ is a geometrically integral component, defined over $\F_q$, of multiplicity 1 and degree $d'$. Denote by $g'$ the geometric genus of the normalization of $C'$. Then we have
    \begin{equation}\label{eq:ns_pt_bound}
        \#C(\F_q)^{\sm} \geq q+1 - g' \lfloor 2\sqrt{q} \rfloor - (d'-1)(d'-2) + 2g' - (d-d')d'.
    \end{equation}
\end{lemma}

\begin{lemma}
\label{lem:plane_with_3pts}
    Let \qbound{591}. Suppose $f \in \F_q[x_0, x_1, x_2]$ is a nonzero ternary degree 7 form with no linear factors defined over $\F_q$. Suppose further that there exist $u,v,w \in X_f(\F_q)$ such that $\Span{u,v,w} = \P^2$. Then either $X_f(\F_q)^{\sm} \neq \emptyset$ or there exist $y,z \in X_f(\F_q)$ such that $\Span{y,z} \not\subset X_f$ and $X_f \cap \Span{y,z}$ contains a geometric point of multiplicity 1.
\end{lemma}

\begin{lemma}
\label{lem:plane_with_3pts_d=11}
    Let \qbound{7061}. Suppose $f \in \F_q[x_0, x_1, x_2]$ is a nonzero ternary degree 11 form with no linear factors defined over $\F_q$. Suppose further that there exist $u,v,w \in X_f(\F_q)$ such that $\Span{u,v,w} = \P^2$. Then either $X_f(\F_q)^{\sm} \neq \emptyset$ or there exist $y,z \in X_f(\F_q)$ such that $\Span{y,z} \not\subset X_f$ and $X_f \cap \Span{y,z}$ contains a geometric point of multiplicity 1.
\end{lemma}

\begin{proof}
    The possible factorization types of $f$ over $\F_q$ are
    \begin{align*}
        (11), (9,2), (8,3), 
        (7,4), (7,2,2), (7,2^2), 
        (6,5), (6,3,2), \\
        (5,4,2), (5,3,3), (5,3^2), (5,2,2,2), (5,2^2,2), (5,2^3),
        (4,4,3), (4^2,3), (4,3,2,2), (4,3,2^2),\\
        (3,3,3,2), (3^2,3,2), (3^3,2), (3,2,2,2,2), (3,2^2,2,2), (3,2^2,2^2), (3,2^3,2), \text{ or } (3,2^4).
    \end{align*}
    If any of the factors of multiplicity 1 defined over $\F_q$ are irreducible over $\overline{\F_q}$, then we can use Lemma \ref{lem:finding_sm_pt} to find $X_f(\F_q)^\sm \neq \emptyset$, as in the proof of Lemma \ref{lem:plane_with_3pts}. The limiting case is that of $f$ geometrically irreducible, where the presence of the three singular $\F_q$-points $u,v,w$ means \qbound{7061} suffices.

    Write the factorization over $\F_q$ as $f = \prod_{1 \leq i \leq t} g_i^{e_i}$, i.e.\ if $f$ has type $(d_1^{e_1},\ldots,d_t^{e_t})$ over $\F_q$, then $g_i$ is defined over $\F_q$ and has degree $d_i$. Assume for all $i$ with $e_i = 1$, we have $g_i$ is geometrically reducible; otherwise we are in the situation above.
    
    Suppose there exists an $i$ with the following properties:
    \begin{enumerate}[label = (\roman*)]
        \item $d_i$ is prime,
        \item $e_i = 1$,
        \item $d_i > \frac{d_j}{2}$ for all $j \neq i$ with $e_j = 1$,
        \item $d_i > d_j$ for all $j \neq i$ with $e_j > 1$.
    \end{enumerate}
    Since $g_i$ is reducible over $\overline{\F_q}$, it must factor into $d_i$ linear forms defined over $\F_{q^{d_i}}$, cyclically permuted by the action of $\mathrm{Gal}(\F_{q^{d_i}}/\F_q)$. In particular, $g_i$ vanishes on at most one $\F_q$-point, so without loss of generality it vanishes on $d_i$ distinct points on the line $\Span{u,v}$ defined (and conjugate) over $\F_{q^{d_i}}$.

    Consider now how the $g_j$ vanish on $\Span{u,v}$ for $j \neq i$. If $e_j = 1$, then $g_j$ --- and its restriction to the line $\Span{u,v}$ --- factors into forms of degree at most $d_j/2$ over $\overline{\F_q}$. By (iii), the degree of these factors of $g_j$ is too small to vanish on any $\F_{q^{d_i}}$-points of the line. The same argument works for $e_j > 1$, since we have the stronger hypothesis (iv). Thus $u,v \in X_f(\F_q)$ satisfy the desired requirements, with the zeros of $g_i$ on this line providing the necessary geometric point(s) of multiplicity 1. 

    Conditions (i) -- (iv) are satisfied for all types except $(11)$, $(9,2)$, $(8,3)$, $(6,3,2)$, $(4^2,3)$, $(3^2,3,2)$, and $(3^3,2)$. We deal with the stragglers ad hoc.

    \begin{itemize}\setlength{\itemsep}{1ex}
        \item[$(11)$] If $f$ is geometrically reducible, then it must factor into 11 conjugate linear forms defined over $\F_{q^{11}}$. However, this is not possible, as such $f$ vanishes on at most one $\F_q$-point.

        \item[$(9,2)$] If $g_1$ is not geometrically irreducible, it must factor into three conjugate cubic forms $h,h',h''$ over $\F_{q^3}$. As above, we may assume the quadratic factor $g_2$ vanishes on two distinct $\F_{q^2}$-points of $\Span{u,v}$. Since $g_1$ must then vanish on both $u,v \in X_f(\F_q)$, each of $h,h',h''$ must vanish on both $u$ and $v$, and hence at most one more $\F_q$-point on $\Span{u,v}$, so the line $\Span{u,v}$ satisfies the desired requirements.

        \item[$(8,3)$] If $g_1$ is not geometrically irreducible, it must factor into two conjugate quartic forms $h,h'$ over $\F_{q^2}$, while we may assume the cubic factor $g_2$ vanishes on three distinct conjugate $\F_{q^3}$-points of $\Span{u,v}$. As in the $(9,2)$ case, $h,h'$ are forced to vanish on $u,v$, so they cannot vanish on any $\F_{q^3}$-points along $\Span{u,v}$.

        \item[$(6,3,2)$] Here $g_1$ can factor in two ways. Suppose it factors into conjugate cubic forms $h,h'$ defined over $\F_{q^2}$ and assume the cubic factor $g_2$ vanishes on three distinct $\F_{q^3}$-points of $\Span{u,v}$. The quadratic factor $g_3$ vanishes on at most one of $u$ or $v$, thus $h,h'$ must both vanish on (at least) the other. Their degrees are then too small to also vanish on any $\F_{q^6}$-points on $\Span{u,v}$, nor in particular any $\F_{q^3}$-points, so we are done.

        If instead $g_1$ factors into conjugate quadratic forms $h,h',h''$ defined over $\F_{q^3}$, a similar argument shows they all must vanish on at least one of $u$ or $v$ and one other $\F_q$-point on the line $\Span{u,v}$, while the quadratic factor $g_3$ vanishes on distinct $\F_{q^2}$-points of $\Span{u,v}$. The cubic factor $g_2$ meets the line at an $\F_q$-point or distinct $\F_{q^3}$-points, but eitherway the requirements of the lemma are met.

        \item[$(4^2,3)$] Assume the cubic factor $g_2$ meets the line $\Span{u,v}$ at three distinct $\F_{q^3}$-points. For degree reasons only, the quartic factor $g_1$ can vanish on at most two of them, since it must vanish on $u$ and $v$.
        
        \item[$(3^2,3,2)$] Suppose the cubic factor $g_1$ is irreducible over $\overline{\F_q}$. The cubic factor $g_2$ vanishes on at most one $\F_q$-point in the plane. Let $y,z \in X_{g_1}(\F_q)$ be distinct points such that $\Span{y,z}$ avoids this point (if it exists); note that the hypotheses on $q$ ensure more than enough $\F_q$-solutions to $g_1=0$ for this. On the line $\Span{y,z}$, $g_1$ vanishes on $y$, $z$, and at least one more point, necessarily defined over $\F_q$; the cubic factor $g_2$ vanishes in three distinct $\F_{q^3}$-points, while the quadratic $g_3$ vanishes on an $\F_q$-point or two $\F_{q^2}$-points. In either case, the zeros of $g_2$ along this line provide the necessary points of multiplicity 1. 

        If $g_1$ is reducible over $\overline{\F_q}$, then in fact $X_f(\F_q) = \{u,v,w\}$, with each of $g_1, g_2, g_3$ vanishing on exactly one. If $u \in X_{g_1}$, then either of $\Span{u,v}$ or $\Span{u,w}$ satisfy the desired requirements.

        \item[$(3^3,2)$] Assume the quadratic factor $g_2$ meets the line $\Span{u,v}$ at two distinct $\F_{q^2}$-points. As in the $(4^2,3)$ case, the cubic factor $g_1$ can vanish on at most one of them, since it must vanish on $u$ and $v$.
\end{itemize}\vspace{-1ex}
\end{proof}

\end{document}
